\documentclass[../../main.tex]{subfiles}

\begin{document}

\section{Riesz-Schauder 定理: Fredholm 算子的刻画}\label{section:ch16-riesz-schauder}

Banach 空间 $X$ 的子空间 $X_{0}$ 称为是在 $X$ 有有限余维的, 若它在 $X$ 有有限维的线性补, 即存在 $X$ 的有限维子空间 $X_{1}$ 使得 $X=X_{0}\oplus X_{1}$. $X_{0}$ 的余维, 记为 $\operatorname{codim}X_{0}$, 恰当地定义为 $X_{0}$ 的线性补的维数; 所有线性补有同样的维数(见习题\cref{exercise:ch16-66}). 线性代数的基石是断言: 若 $X$ 是有限维线性空间而 $T:X\to X$ 是线性的, 则 $T$ 的秩与 $T$ 的零空间的维数的和等于 $X$ 的维数, 即若 $\dim X=n$,
\begin{align*}
\dim\operatorname{Im}T+\dim\ker T=n.
\end{align*}
因此, 由于 $\operatorname{codim}\operatorname{Im}T=n-\dim\operatorname{Im}T$,
\begin{align}
\dim\ker T=\operatorname{codim}\operatorname{Im}T.\label{align:ch16-fredholm-basic}
\end{align}

本节我们的主要目标是证明: 若 $H$ 是 Hilbert 空间, 而算子 $T\in\mathcal{L}(H)$ 是恒等算子的紧扰动, 则 $T$ 有使得\eqref{align:ch16-fredholm-basic}成立的有限维的核与有限余维的象.

\begin{proposition}\label{proposition:ch16-idplusK-kernel}
令 $H$ 为 Hilbert 空间而 $K\in\mathcal{L}(H)$ 是紧的, 则 $\operatorname{Id}+K$ 有有限维的核与闭的象.
\end{proposition}

\begin{proof}
假定 $\ker(\operatorname{Id}+K)$ 是无穷维的. 我们从定理\cref{theorem:ch16-separable-ono-basis}推出, 存在包含于 $\ker(\operatorname{Id}+K)$ 的单位向量的正交序列 $\{u_{k}\}$. 由于若 $m\neq n$, $\|K(u_{n})-K(u_{m})\|=\|u_{n}-u_{m}\|=\sqrt{2}$, 序列 $\{K(u_{n})\}$ 没有收敛的子序列. 这与算子 $K$ 的紧性矛盾. 因此 $\dim[\ker(\operatorname{Id}+K)]<\infty$. 令 $H_{0}=[\ker(\operatorname{Id}+K)]^{\perp}$. 我们宣称存在 $c>0$ 使得
\begin{align}
\text{对所有 }u\in H_{0},\quad\|u+K(u)\|\geqslant c\|u\|.\label{align:ch16-idplusK-est}
\end{align}
事实上, 若不存在这样的 $c$, 则选取 $H_{0}$ 中的单位向量序列 $\{u_{n}\}$, 使得在 $H$ 中弱意义下 $\{u_{n}+K(u_{n})\}\to 0$. 由于 $K$ 是紧的, 若必要的话过渡到子序列, 我们可以假定在强意义下 $\{K(u_{n})\}\to h_{0}$. 因此在强意义下 $\{u_{n}\}\to-h_{0}$. 根据 $K$ 的连续性, $h_{0}+K(h_{0})=0$. 因此, 由于 $H_{0}$ 是闭的, $h_{0}$ 是既属于 $[\ker(\operatorname{Id}+K)]^{\perp}$ 又属于 $\ker(\operatorname{Id}+K)$ 的单位向量. 这个矛盾证明了使得\eqref{align:ch16-idplusK-est}成立的 $c$ 的存在性. 我们从\eqref{align:ch16-idplusK-est}与 $H_{0}$ 的完备性推出 $(\operatorname{Id}+K)(H_{0})$ 是闭的. 由于 $(\operatorname{Id}+K)(H_{0})=(\operatorname{Id}+K)(H)$, $\operatorname{Im}(\operatorname{Id}+K)$ 是闭的.
\end{proof}

\begin{definition}[正交投影序列]\label{definition:ch16-proj-sequence}
令 $\{\varphi_{n}\}$ 为可分 Hilbert 空间 $H$ 的规范正交基. 对每个 $n$, 定义 $P_{n}\in\mathcal{L}(H)$ 为
\begin{align*}
\text{对所有 }h\in H,\quad P_{n}(h)=\sum_{k=1}^{n}\langle\varphi_{k},h\rangle\varphi_{k}.
\end{align*}
我们称 $\{P_{n}\}$ 为由 $\{\varphi_{n}\}$ 诱导的正交投影序列.
\end{definition}

对由规范正交基 $\{\varphi_{n}\}$ 诱导的正交投影序列 $\{P_{n}\}$, 每个 $P_{n}$ 是从 $H$ 到 $\operatorname{span}\{\varphi_{1},\ \cdots,\ \varphi_{n}\}$ 的正交投影, 因此 $\|P_{n}\|=1$. 此外, 根据规范正交基的定义, 在 $H$ 上逐点地 $\{P_{n}\}\to\operatorname{Id}$. 因此, 对任何 $T\in\mathcal{L}(H)$, $\{P_{n}\circ T\}$ 是在 $H$ 上逐点收敛于 $T$ 的有限秩算子序列.

\begin{proposition}\label{proposition:ch16-finite-rank-approx}
令 $\{P_{n}\}$ 为可分 Hilbert 空间的规范正交基 $\{\varphi_{n}\}$ 诱导的正交投影序列. 则算子 $T\in\mathcal{L}(H)$ 是紧的当且仅当在 $\mathcal{L}(H)$ 中 $\{P_{n}\circ T\}\to T$.
\end{proposition}

\begin{proof}
首先假设在 $\mathcal{L}(H)$ 中 $\{P_{n}\circ T\}\to T$. 对每个自然数 $n$, $P_{n}\circ T$ 有有限维的象, 因此 $(P_{n}\circ T)(B)$ 是全有界的, 其中 $B$ 是 $H$ 的单位球. 由于 $\{P_{n}\circ T\}\to T\in\mathcal{L}(H)$, 在 $B$ 上 $\{P_{n}\circ T\}$ 一致收敛于 $T$. 因此 $T(B)$ 也是全有界的. 我们从命题\cref{proposition:ch16-compact-equiv}推出算子 $T$ 是紧的. 为证明逆命题, 假设 $T$ 是紧的. 则集合 $\overline{T(B)}$ 关于强拓扑是紧的. 对每个自然数 $n$, 定义 $\psi_{n}:\overline{T(B)}\to\mathbf{R}$ 为
\begin{align*}
\text{对所有 }h\in\overline{T(B)},\quad\psi_{n}(h)=\|P_{n}(h)-h\|.
\end{align*}
由于每个 $P_{n}$ 范数为 $1$, 实值函数序列 $\{\psi_{n}:\overline{T(B)}\to\mathbf{R}\}$ 是等度连续的、有界的, 且在紧集 $\overline{T(B)}$ 上逐点收敛到 $0$. 我们从 Arzel\`a-Ascoli 定理推出 $\{\psi_{n}:\overline{T(B)}\to\mathbf{R}\}$ 一致收敛到 $0$. 这正好意味着在 $\mathcal{L}(H)$ 中 $\{P_{n}\circ T\}\to T$.
\end{proof}

\begin{proposition}\label{proposition:ch16-idplusK-onto}
令 $H$ 为 Hilbert 空间而 $K\in\mathcal{L}(H)$ 是紧的. 若 $\operatorname{Id}+K$ 是一对一的, 则它是映上的.
\end{proposition}

\begin{proof}
我们把证明存在 $H$ 的闭可分子空间 $H_{0}$ 使得 $K(H_{0})\subseteq H_{0}$ 且在 $H_{0}^{\perp}$ 上 $K=0$ 留作练习(习题\cref{exercise:ch16-53}), 通过用 $H_{0}$ 代替 $H$ 我们可以假定 $H$ 是可分的. 由于 $\operatorname{Id}+K$ 是一对一的, 我们看到 $\ker(\operatorname{Id}+K)^{\perp}=H$. 利用这个事实, 如同我们在命题\cref{proposition:ch16-idplusK-kernel}的证明中所做的那样, 我们可证明存在 $c>0$ 使得
\begin{align}
\text{对所有 }h\in H,\quad\|h+K(h)\|\geqslant c\|h\|.\label{align:ch16-idplusK-onto-1}
\end{align}
根据定理\cref{theorem:ch16-separable-ono-basis}, $H$ 有规范正交基 $\{\varphi_{n}\}$. 令 $\{P_{n}\}$ 为由 $\{\varphi_{n}\}$ 诱导的正交投影序列. 对每个自然数 $n$, 令 $H_{n}$ 为 $\{\varphi_{1},\ \cdots,\ \varphi_{n}\}$ 的线性扩张. 由于算子 $K$ 是紧的, 根据前一个命题, 在 $\mathcal{L}(H)$ 中 $\{P_{n}\circ K\}\to K$. 选取自然数 $N$ 使得对所有 $n\geqslant N$, $\|P_{n}\circ K-K\|<c/2$. 我们从\eqref{align:ch16-idplusK-onto-1}推出
\begin{align}
\text{对所有 }u\in H\text{ 与所有 }n\geqslant N,\quad\|u+P_{n}\circ K(u)\|\geqslant c/2\|u\|.\label{align:ch16-idplusK-onto-2}
\end{align}
为证明 $(\operatorname{Id}+K)(H)=H$, 令 $h_{*}$ 属于 $H$. 令 $n\geqslant N$. 从\eqref{align:ch16-idplusK-onto-2}得出 $\operatorname{Id}+P_{n}\circ K$ 在 $H_{n}$ 上的限制是一对一线性算子. 它将有限维空间 $H_{n}$ 映到其自身. 有限维空间上的一对一线性算子是映上的. 因此该限制将 $H_{n}$ 映上 $H_{n}$. 因此, 存在向量 $u_{n}\in H_{n}$ 使得
\begin{align}
u_{n}+(P_{n}\circ K)(u_{n})=P_{n}(h_{*}).\label{align:ch16-idplusK-onto-3}
\end{align}
该等式两边取关于 $v\in H$ 的内积且用投影 $P_{n}$ 的对称性得出
\begin{align}
\text{对所有 }n\geqslant N,v\in H,\quad\langle u_{n}+K(u_{n}),P_{n}(v)\rangle=\langle h_{*},P_{n}(v)\rangle.\label{align:ch16-idplusK-onto-4}
\end{align}
我们从\eqref{align:ch16-idplusK-onto-3}与估计式\eqref{align:ch16-idplusK-onto-2}推出
\begin{align*}
\text{对所有 }n\geqslant N,\quad\|h_{*}\|\geqslant\|P_{n}(h_{*})\|=\|u_{n}+(P_{n}\circ K)u_{n}\|\geqslant c/2\|u_{n}\|.
\end{align*}
因此序列 $\{u_{n}\}$ 是有界的. 定理\cref{theorem:ch16-weak-seq}告诉我们存在弱收敛于 $u\in H$ 的子序列 $\{u_{n_{k}}\}$. 因此 $\{u_{n_{k}}+K(u_{n_{k}})\}$ 弱收敛于 $u+K(u)$. 在\eqref{align:ch16-idplusK-onto-4}取 $k\to\infty$ 的极限, 由命题\cref{proposition:ch16-weak-properties}得出
\begin{align*}
\text{对所有 }v\in H,\quad\langle u+K(u),v\rangle=\langle h_{*},v\rangle.
\end{align*}
因此 $u+K(u)=h_{*}$. 因此 $(\operatorname{Id}+K)(H)=H$.
\end{proof}

\begin{theorem}[Riesz-Schauder]\label{theorem:ch16-riesz-schauder}
令 $H$ 为 Hilbert 空间而 $K\in\mathcal{L}(H)$ 是紧的. 则 $\operatorname{Im}(\operatorname{Id}+K)$ 是闭的且
\begin{align}
\dim\ker(\operatorname{Id}+K)=\dim\ker(\operatorname{Id}+K^{*})<\infty.\label{align:ch16-riesz-schauder-1}
\end{align}
特别地, $\operatorname{Id}+K$ 是一对一的当且仅当它是映上的.
\end{theorem}

\begin{proof}
根据命题\cref{proposition:ch16-idplusK-kernel}, 恒等算子的紧扰动有有限维的核与闭的象. 我们将证明
\begin{align}
\dim\ker(\operatorname{Id}+K)\geqslant\dim\ker(\operatorname{Id}+K^{*}).\label{align:ch16-riesz-schauder-2}
\end{align}
一旦这得到证明, 我们用 $K^{*}$ 代替 $K$ 且观察到 $(K^{*})^{*}=K$, 以及关于 $K^{*}$ 的紧性的 Schauder 定理, 得到反向的不等式. 我们用反证法证明\eqref{align:ch16-riesz-schauder-2}. 若不是这样, $\dim\ker(\operatorname{Id}+K)<\dim\ker(\operatorname{Id}+K^{*})$. 令 $P$ 为 $H$ 映上 $\ker(\operatorname{Id}+K)$ 的正交投影, 而 $A$ 是 $\ker(\operatorname{Id}+K)$ 到 $\ker(\operatorname{Id}+K^{*})$ 的一对一但不是映上的线性映射. 定义 $K^{\prime}=K+A\circ P$. 由于 $\operatorname{Id}+K$ 有闭的象, 命题\cref{proposition:ch16-kernel-orth}告诉我们,
\begin{align*}
H=\operatorname{Im}(\operatorname{Id}+K)\oplus\ker(\operatorname{Id}+K^{*}).
\end{align*}
因此 $\operatorname{Id}+K^{\prime}$ 是一对一但不是映上的. 另一方面, 由于 $A\circ P$ 是有限秩的, 它是紧的, 因此 $K^{\prime}$ 也如此. 这两个断言与前一个命题矛盾. 因此\eqref{align:ch16-riesz-schauder-2}得证. 由于 $\operatorname{Id}+K$ 有闭的象, 我们从\eqref{align:ch16-coercive}与\eqref{align:ch16-riesz-schauder-2}推出 $\operatorname{Id}+K$ 是一对一的当且仅当它是映上的.
\end{proof}

\begin{corollary}[Fredholm 择一定理]\label{corollary:ch16-fredholm-alternative}
令 $H$ 为 Hilbert 空间, $K\in\mathcal{L}(H)$ 是紧的而 $\mu$ 是非零实数. 则以下结论恰有一个成立:
\begin{enumerate}[(i)]
\item 方程
\begin{align*}
\mu h-K(h)=0,\ h\in H
\end{align*}
有非零解.
\item 对每个 $h_{0}\in H$, 方程
\begin{align*}
\mu h-K(h)=h_{0},\ h\in H
\end{align*}
有唯一解.
\end{enumerate}
\end{corollary}

\begin{definition}[Fredholm 算子, 指标]\label{definition:ch16-fredholm}
令 $H$ 为 Hilbert 空间而 $T$ 属于 $\mathcal{L}(H)$. 则 $T$ 称为是 Fredholm 的, 若 $T$ 的核是有限维的且 $T$ 的象有有限余维. 对这样的算子, 它的指标 $\operatorname{ind}T$ 定义为
\begin{align*}
\operatorname{ind}T=\dim\ker T-\operatorname{codim}\operatorname{Im}T.
\end{align*}
\end{definition}

在 Riesz-Schauder 定理的证明中, 我们首先证明 $\operatorname{Im}T$ 是闭的且用此以及\eqref{align:ch16-kernel-orth}来证明 $\operatorname{Im}T$ 有等于 $\dim\ker T^{*}$ 的有限余维. 然而, 第 13 章的定理 12 告诉我们, 若 $H$ 是 Hilbert 空间而算子 $T\in\mathcal{L}(H)$ 有有限余维的象, 则它的象是闭的. 因此每个 Fredholm 算子有闭的象, 又一次由\eqref{align:ch16-kernel-orth}, $\operatorname{codim}\operatorname{Im}T=\dim\ker T^{*}$.

我们说算子 $T$ 是可逆的, 若它是一对一与映上的. 开映射定理告诉我们, 可逆算子的逆是连续的, 所以可逆算子是一个同构.

\begin{theorem}\label{theorem:ch16-index-zero}
令 $H$ 为 Hilbert 空间而 $T$ 属于 $\mathcal{L}(H)$. 则 $T$ 是指标为 $0$ 的 Fredholm 算子当且仅当 $T=S+K$, 其中 $S\in\mathcal{L}(H)$ 是可逆的而 $K\in\mathcal{L}(H)$ 是紧的.
\end{theorem}

\begin{proof}
首先假设 $T$ 是指标为 $0$ 的 Fredholm 算子. 由于 $\operatorname{Im}T$ 是闭的, 命题\cref{proposition:ch16-kernel-orth}告诉我们,
\begin{align}
H=\operatorname{Im}T\oplus\ker T^{*}.\label{align:ch16-index-zero-1}
\end{align}
由于 $\dim\ker T=\dim\ker T^{*}<\infty$, 我们可以选取将 $\ker T$ 映上 $\ker T^{*}$ 的一对一算子 $A$. 令 $P$ 为 $H$ 映上 $\ker T$ 的正交投影算子. 定义 $K=A\circ P\in\mathcal{L}(H)$ 与 $S=T-K$. 则 $T=S+K$. 算子 $K$ 是紧的, 由于它是有限秩的, 而根据\eqref{align:ch16-index-zero-1}与 $P$ 和 $A$ 的选取, 算子 $S$ 是可逆的. 因此 $T$ 是可逆算子的紧扰动.

为证明逆命题, 假定 $T=S+K$, 其中 $S\in\mathcal{L}(H)$ 是可逆的而 $K\in\mathcal{L}(H)$ 是紧的. 观察到
\begin{align}
T=S\circ[\operatorname{Id}+S^{-1}\circ K].\label{align:ch16-index-zero-2}
\end{align}
由于 $S^{-1}$ 是连续的且 $K$ 是紧的, $S^{-1}\circ K$ 是紧的. Riesz-Schauder 定理告诉我们 $\operatorname{Id}+S^{-1}\circ K$ 是指标为 $0$ 的 Fredholm 算子. Fredholm 算子与可逆算子的复合也是指标为 $0$ 的 Fredholm 算子(见习题\cref{exercise:ch16-55}). 我们因此从\eqref{align:ch16-index-zero-2}推出 $T$ 是指标为 $0$ 的 Fredholm 算子.
\end{proof}

我们把下面推论的证明留作练习.

\begin{corollary}\label{corollary:ch16-fredholm-composition}
令 $H$ 为 Hilbert 空间, 而 $T$ 与 $S$ 是 $\mathcal{L}(H)$ 中指标为 $0$ 的 Fredholm 算子. 则复合 $S\circ T$ 也是指标为 $0$ 的 Fredholm 算子.
\end{corollary}

\medskip\par\noindent\textbf{注. }Riesz-Schauder 定理和定理\cref{theorem:ch16-index-zero}对一般的 Banach 空间上的算子也成立, 其中紧线性算子定义为将单位球映到有紧闭包的集合的算子. 然而, 证明的一般方法十分不同. 命题\cref{proposition:ch16-finite-rank-approx}的证明要点在于用有限秩的线性算子逼近 $\mathcal{L}(H)$ 中的紧线性算子. Per Enflo 已证明了可分的 Banach 空间上存在不能被有限秩线性算子在 $\mathcal{L}(H)$ 中逼近的线性紧算子.

\begin{exercise}\label{exercise:ch16-53}
令 $K\in\mathcal{L}(H)$ 为紧算子. 证明 $T=K^{*}K$ 是紧与对称的. 接着用 Hilbert-Schmidt 定理证明: 存在 $H$ 的规范正交序列 $\{\varphi_{k}\}$ 使得对所有 $k$, $T(\lambda_{k}\varphi_{k})=\lambda_{k}\varphi_{k}$, 且若 $h$ 正交于 $\{\varphi_{k}\}_{k=0}^{\infty}$, 则 $T(h)=0$. 得出若 $h$ 正交于 $\{\varphi_{k}\}_{k=0}^{\infty}$, 则
\begin{align*}
\|K(h)\|^{2}=\langle K(h),K(h)\rangle=\langle T(h),h\rangle=0.
\end{align*}
定义 $H_{0}$ 为 $\{K^{m}(\varphi_{k})\mid m\geqslant 0,\ k\geqslant 1\}$ 的闭线性扩张. 证明 $H_{0}$ 是闭与可分的, $K(H_{0})\subseteq H_{0}$ 且在 $H_{0}^{\perp}$ 上 $K=0$.
\end{exercise}

\begin{exercise}\label{exercise:ch16-54}
令 $\mathcal{K}(H)$ 表示 $\mathcal{L}(H)$ 中的紧算子集. 证明 $\mathcal{K}(H)$ 是 $\mathcal{L}(H)$ 的闭子空间, 它有有限秩算子集作为它的稠密子空间. $\mathcal{K}(H)$ 是 $\mathcal{L}(H)$ 的开子集吗?
\end{exercise}

\begin{exercise}\label{exercise:ch16-55}
证明指标为 $0$ 的 Fredholm 算子与可逆算子以两种次序的复合也是指标为 $0$ 的 Fredholm 算子.
\end{exercise}

\begin{exercise}\label{exercise:ch16-56}
证明两个指标为 $0$ 的 Fredholm 算子的复合是指标为 $0$ 的 Fredholm 算子.
\end{exercise}

\begin{exercise}\label{exercise:ch16-57}
证明算子 $T\in\mathcal{L}(H)$ 是指标为 $0$ 的 Fredholm 算子当且仅当它是可逆算子通过有限秩算子的扰动.
\end{exercise}

\begin{exercise}\label{exercise:ch16-58}
用以下方法证明 $\mathcal{L}(H)$ 中的可逆算子全体是 $\mathcal{L}(H)$ 的开子集.
\begin{enumerate}[(i)]
\item 对满足 $\|A\|<1$ 的 $A\in\mathcal{L}(H)$, 用 $\mathcal{L}(H)$ 的完备性证明 Neumann 级数 $\sum_{n=0}^{\infty}A^{n}$ 收敛于 $\mathcal{L}(H)$ 中的算子, 该算子是 $\operatorname{Id}-A$ 的逆.
\item 对可逆算子 $S\in\mathcal{L}(H)$, 证明对任何 $T\in\mathcal{L}(H)$, $T=S[\operatorname{Id}+S^{-1}(T-S)]$.
\item 用(i)和(ii)证明: 若 $S\in\mathcal{L}(H)$ 是可逆的, 则任何满足 $\|S-T\|<1/\|S^{-1}\|$ 的 $T\in\mathcal{L}(H)$ 也是可逆的.
\end{enumerate}
\end{exercise}

\begin{exercise}\label{exercise:ch16-59}
证明指标为 $0$ 的 Fredholm 算子的集合是 $\mathcal{L}(H)$ 的开子集.
\end{exercise}

\begin{exercise}\label{exercise:ch16-60}
用在命题\cref{proposition:ch16-idplusK-onto}的证明中的正交逼近序列方法给出当 $H$ 是可分的情形时命题\cref{proposition:ch16-coercive-invertible}的另一个证明.
\end{exercise}

\begin{exercise}\label{exercise:ch16-61}
对 $T\in\mathcal{L}(H)$, 假定对所有 $h\in H$, $\langle T(h),\ h\rangle\geqslant\|h\|^{2}$. 假设 $K\in\mathcal{L}(H)$ 是紧的而 $T+K$ 是一对一的. 证明 $T+K$ 是映上的.
\end{exercise}

\begin{exercise}\label{exercise:ch16-62}
令 $K\in\mathcal{L}(H)$ 是紧的而 $\mu\in\mathbf{R}$ 且 $|\mu|>\|K\|$. 证明 $\mu-K$ 是可逆的.
\end{exercise}

\begin{exercise}\label{exercise:ch16-63}
令 $S\in\mathcal{L}(H)$ 满足 $\|S\|<1$, $K\in\mathcal{L}(H)$ 是紧的且 $(\operatorname{Id}+S+K)(H)=H$. 证明 $\operatorname{Id}+S+K$ 是一对一的.
\end{exercise}

\begin{exercise}\label{exercise:ch16-64}
令 $\mathcal{G}\mathcal{L}(H)$ 表示 $\mathcal{L}(H)$ 中的可逆算子集.
\begin{enumerate}[(i)]
\item 证明在算子的复合运算下, $\mathcal{G}\mathcal{L}(H)$ 是一个群: 它称为 $H$ 的一般线性群.
\item $\mathcal{G}\mathcal{L}(H)$ 中的算子 $T$ 称为是正交的, 若 $T^{*}=T^{-1}$. 证明正交算子集是 $\mathcal{G}\mathcal{L}(H)$ 的一个子群: 称之为正交群.
\end{enumerate}
\end{exercise}

\begin{exercise}\label{exercise:ch16-65}
令 $H$ 为 Hilbert 空间, $T\in\mathcal{G}\mathcal{L}(H)$ 是指标为 $0$ 的 Fredholm 算子, 而 $K\in\mathcal{L}(H)$ 是紧的. 证明 $T+K$ 是指标为 $0$ 的 Fredholm 算子.
\end{exercise}

\begin{exercise}\label{exercise:ch16-66}
令 $X_{0}$ 为 Banach 空间 $X$ 的有限余维子空间. 证明 $X_{0}$ 在 $X$ 中的所有有限维线性补有相同的维数.
\end{exercise}

\end{document}